\begin{tikzpicture}[font=\small, >=latex, node distance=7mm and 9mm]
  \tikzset{flowbox/.style={draw, rounded corners=2pt, align=center, minimum height=11mm, text width=29mm},
    input/.style={flowbox, fill=orange!12, draw=orange!70!black},
    check/.style={flowbox, fill=green!10, draw=green!55!black},
    run/.style={flowbox, fill=blue!10, draw=blue!65!black},
    output/.style={flowbox, fill=purple!10, draw=purple!65!black}}
  \node[input] (csv) {Seven-column\\{\bfseries student CSV}};
  \node[check, right=of csv] (schema) {Schema and\\identifier checks};
  \node[run, right=of schema] (model) {Released planning\\context + validation\\inputs};
  \node[output, right=of model, yshift=12mm] (preview) {Student preview\\six-hour sampled,\\idle-storage,\\non-authoritative screening};
  \node[output, right=of model, yshift=-12mm] (eval) {Exact instructor evaluator\\joint storage-aware\\DCOPF/LP\\authoritative result};
  \draw[->, thick] (csv) -- (schema);
  \draw[->, thick] (schema) -- (model);
  \draw[->, thick] (model) -- (preview);
  \draw[->, thick] (model) -- (eval);
  \node[below=4mm of schema, align=center, text width=42mm] {Reject stale fields and malformed units before any model run.};
  \node[below=4mm of model, align=center, text width=44mm] {The released ACTIVSg2000 package supplies the context and validation inputs; the CSV supplies the student's site capacities.};
\end{tikzpicture}
